\section{Discussion}
\label{sec:discussion}

\paragraph{The mesh was the easy part.} We expected generated meshes to fail as
meshes: holes, non-manifold edges, stray fragments. They did not. The problems
that cost development time were conventions that a file does not state (size,
facing), a texture budget set by the generator rather than the user, a processing
step whose input representation we did not check, and an animation whose pose
differed from its bind pose. Three of the six fixes repaired our own processing,
not the generator's output, and the one simplification step we added both tore
the mesh and rested on a wrong diagnosis: its ``UV floor'' was a loop that ran
out of rounds. For a developer
adopting such a pipeline, the lesson is less ``inspect the generator'' than
``inspect every step you add after it, and the explanation you wrote down for it''.

\paragraph{A pre-import checklist.} Ordered by what each check would have caught
in this game, at the cost of a script that runs in seconds per asset:
\begin{enumerate}
\item \emph{Size and facing.} Read the largest extent and treat it as unitless;
  supply real size from the design. Facing is not in the file: render one
  thumbnail from $+Z$, or adopt the generator's convention once and test it on
  the first asset.
\item \emph{Texture budget.} Count texels, not triangles. Three 4096$^2$ maps cost
  64\,MiB in BC7 per asset; most of the half-empty atlas can be dropped by
  downscaling or re-baking.
\item \emph{Simplification input.} Weld before simplifying, and check open edges
  after. A seam-split buffer looks like a mesh and simplifies without error.
\item \emph{Animated contact.} For rigged assets, evaluate each loop offline and
  compare its lowest vertex with the bind pose; budget the bob amplitude, or use
  foot contact handling.
\item \emph{Topology.} Watertightness and manifoldness were never the problem
  here, but they are cheap to count and our own tool was the one that broke them.
\end{enumerate}

\paragraph{Counting versus looking.} Our earlier audit of an LLM-generated quiz
bank found that the defects that matter were invisible one item at a
time and obvious when counted across the set~\cite{zhao2026countingwhat}. Generated
meshes partly reverse this. Counting confirmed that every asset shares the same
conventions, so one correction serves all of them, and counting found the torn
companion and the texture budget. But the costliest defects, facing and grounding,
were found by looking at the game, and could only have been found earlier by
looking: rendering a view, or evaluating an animation. A useful pre-import check
for generated geometry is a small render and a short simulation, not only a
statistic.
